\section{Screening order and within-agent ranking}
We held the prediction bank, risk screen, and query edit counts fixed on
72 repeated views of 24 fitting recordings from nine source groups.
The nominate-first arm passes \NominationNominateFirst{}/72 recording
views, retaining OLS4 in addition passes \NominationAddBack{}/72, and
screening the full menu first passes \NominationFullScreen{}/72.
Pair coupling after full screening passes \NominationJoint{}/72.
None of the full-screen calibration configurations qualify.

We next refit two matched score heads per cell. Both regress paired
CV-anchored bounds, with candidate displacement from CV and OLS4 in the
input. The second also receives a weighted logistic preference loss
between two paths for the same agent. A preference receives weight only
when its advantage on known steps exceeds the sum of candidate
separations on missing steps, so the signed preference holds for every
completion of those missing positions. No preference is assigned to
unresolved pairs. This objective changes the ranking, not risk
eligibility or the interpretation of a missing reference.

At 2,317 edits per arm, the pointwise and preference models pass
\NominationPoint{}/72 and \NominationPaired{}/72. The direct conservative
utility contrast is \NominationContrast{} percentage points (3,000
source-group draws after averaging seeds). Worst known signed easy
degradation falls from \NominationPointEasy{}\% to
\NominationPairedEasy{}\%, but easy support is undefined in
\NominationUndefinedEasy{} views and the all-view bound exceeds 2\%
in \NominationAllOver{} views. Of \NominationSharedChanged{}
path changes where both select the same recipient, preference gives
\NominationSharedWins{} definite wins and \NominationSharedLosses{}
definite losses; these are not independent trials. Calibration qualifies
\NominationPointQual{}/9 and \NominationPairedQual{}/9 respectively.
This controlled fitting-source comparison is not independent confirmation.
